\documentclass[12pt,a4paper]{article}
\usepackage{fontspec}
\usepackage{xeCJK}
\setCJKmainfont{Microsoft JhengHei}
\setCJKsansfont{Microsoft JhengHei}
\setCJKmonofont{Microsoft JhengHei}
\usepackage{amsmath,amssymb,amsthm}
\usepackage{geometry}
\geometry{margin=2cm}
\usepackage{enumitem}
\usepackage{fancyhdr}
\usepackage{tcolorbox}
\tcbuselibrary{theorems,skins,breakable}
\usepackage{hyperref}
\hypersetup{colorlinks=true,linkcolor=blue!70!black}

\pagestyle{fancy}
\fancyhf{}
\fancyhead[L]{\small 複變函數期末模擬試題}
\fancyhead[R]{\small \thepage}
\setlength{\headheight}{14pt}
\renewcommand{\headrulewidth}{0.4pt}

\newtcolorbox{qbox}[1][]{
  colback=blue!5!white, colframe=blue!60!black,
  fonttitle=\bfseries, title=#1, breakable
}
\newtcolorbox{solbox}[1][]{
  colback=green!3!white, colframe=green!50!black,
  fonttitle=\bfseries, title=#1, breakable
}
\newtcolorbox{ansbox}{
  colback=red!5!white, colframe=red!60!black,
  fonttitle=\bfseries, title=答案, breakable
}

\newcommand{\dd}{\,\mathrm{d}}
\newcommand{\ii}{\mathrm{i}}
\newcommand{\ee}{\mathrm{e}}
\newcommand{\Res}{\operatorname{Res}}

\begin{document}

\begin{center}
{\LARGE\bfseries 複變函數期末模擬試題}\\[6pt]
{\large 共 14 題（含詳解）}\\[4pt]
\end{center}

\bigskip

%==================================================================
% Q1
%==================================================================
\begin{qbox}[第 1 題\,---\,ML 不等式]
$C:|z|=5$，用 ML 不等式求 $\displaystyle\left|\oint_C\frac{\ee^{3z}}{z-2}\dd{z}\right|$ 的上界。若結果為 $a\pi\ee^b$，求 $a+b$。
\end{qbox}

\begin{solbox}[解]
$L=2\pi\cdot 5=10\pi$。

$|\ee^{3z}|=\ee^{3\text{Re}(z)}\le\ee^{15}$（$\text{Re}(z)\le 5$）。

$|z-2|\ge||z|-2|=3$，故 $\dfrac{1}{|z-2|}\le\dfrac{1}{3}$。

$M=\dfrac{\ee^{15}}{3}$。

$$\left|\oint_C\frac{\ee^{3z}}{z-2}\dd{z}\right|\le ML=\frac{\ee^{15}}{3}\cdot 10\pi=\frac{10\pi\ee^{15}}{3}$$

$a=\dfrac{10}{3}$，$b=15$。
\end{solbox}

\begin{ansbox}
$a+b=\dfrac{10}{3}+15=\boxed{\dfrac{55}{3}}$
\end{ansbox}

%==================================================================
% Q2
%==================================================================
\begin{qbox}[第 2 題\,---\,Cauchy--Goursat 與部分分式]
計算 $\displaystyle\oint_C\frac{6z}{(z+2)(z-1)(z-4)}\dd{z}$

(a) $C:|z|=3$，求積分值。

(b) $C:|z|=5$，求積分值。
\end{qbox}

\begin{solbox}[解]
部分分式：$\dfrac{6z}{(z+2)(z-1)(z-4)}=\dfrac{A}{z+2}+\dfrac{B}{z-1}+\dfrac{C}{z-4}$

$A=\dfrac{6(-2)}{(-2-1)(-2-4)}=\dfrac{-12}{18}=-\dfrac{2}{3}$

$B=\dfrac{6(1)}{(1+2)(1-4)}=\dfrac{6}{-9}=-\dfrac{2}{3}$

$C=\dfrac{6(4)}{(4+2)(4-1)}=\dfrac{24}{18}=\dfrac{4}{3}$

\textbf{(a) $|z|=3$}：$z=-2$ 和 $z=1$ 在圓內，$z=4$ 在圓外。
$$=\left(-\frac{2}{3}+\left(-\frac{2}{3}\right)\right)\cdot 2\pi\ii=-\frac{4}{3}\cdot 2\pi\ii=-\frac{8\pi\ii}{3}$$

\textbf{(b) $|z|=5$}：三個奇異點都在圓內。
$$=\left(-\frac{2}{3}-\frac{2}{3}+\frac{4}{3}\right)\cdot 2\pi\ii=0$$
\end{solbox}

\begin{ansbox}
(a) $\boxed{-\dfrac{8\pi\ii}{3}}$ \qquad (b) $\boxed{0}$
\end{ansbox}

%==================================================================
% Q3
%==================================================================
\begin{qbox}[第 3 題\,---\,複數線積分]
計算 $\displaystyle\int_C(x-\ii y^2)\dd{z}$，$C$ 為從 $1$ 到 $1+2\ii$ 的直線段。求實部值。
\end{qbox}

\begin{solbox}[解]
$C$：$z(t)=1+2\ii t$，$0\le t\le 1$。故 $x=1$，$y=2t$，$\dd{z}=2\ii\dd{t}$。

$$\int_C(x-\ii y^2)\dd{z}=\int_0^1(1-\ii(2t)^2)\cdot 2\ii\dd{t}=\int_0^1(1-4\ii t^2)\cdot 2\ii\dd{t}$$
$$=\int_0^1(2\ii-8\ii^2 t^2)\dd{t}=\int_0^1(2\ii+8t^2)\dd{t}$$
$$=\left[8\cdot\frac{t^3}{3}+2\ii t\right]_0^1=\frac{8}{3}+2\ii$$
\end{solbox}

\begin{ansbox}
實部 $=\boxed{\dfrac{8}{3}}$
\end{ansbox}

%==================================================================
% Q4
%==================================================================
\begin{qbox}[第 4 題\,---\,Cauchy 不等式]
$C:|z|=3$，$f(z)=\sin z$，用 Cauchy's inequality 求 $|f^{(2)}(0)|$ 的上界。（提示：$\sinh 3\approx 10.018$）
\end{qbox}

\begin{solbox}[解]
$z_0=0$，$R=3$，$n=2$。

在 $|z|=3$ 上：$|\sin z|=\left|\dfrac{\ee^{\ii z}-\ee^{-\ii z}}{2\ii}\right|\le\dfrac{|\ee^{\ii z}|+|\ee^{-\ii z}|}{2}$

$|\ee^{\ii z}|=\ee^{-\text{Im}(z)}$，在 $|z|=3$ 上，$|\text{Im}(z)|\le 3$，故 $|\ee^{\ii z}|\le\ee^3$，$|\ee^{-\ii z}|\le\ee^3$。

$M=\dfrac{\ee^3+\ee^{-3}}{2}=\cosh 3$。但更簡單的上界：$|\sin z|\le\sinh|z|=\sinh 3$。

實際更精確：$|\sin z|\le\cosh(\text{Im}(z))\le\cosh 3$。不過用 $\sinh 3$ 也可。

取 $M=\sinh 3\approx 10.018$：

$$|f''(0)|\le\frac{2!\cdot\sinh 3}{3^2}=\frac{2\sinh 3}{9}\approx\frac{20.036}{9}\approx 2.226$$

（實際上 $f''(0)=-\sin 0=0$，但 Cauchy 不等式給出上界。）
\end{solbox}

\begin{ansbox}
$|f''(0)|$ 的上界 $\le\boxed{\dfrac{2\sinh 3}{9}\approx 2.23}$
\end{ansbox}

%==================================================================
% Q5
%==================================================================
\begin{qbox}[第 5 題\,---\,Maclaurin 展開與收斂半徑]
將 $f(z)=\dfrac{1}{3+z}$ 展開為 Maclaurin series，求收斂半徑 $R$。
\end{qbox}

\begin{solbox}[解]
$$f(z)=\frac{1}{3+z}=\frac{1}{3}\cdot\frac{1}{1+z/3}=\frac{1}{3}\sum_{k=0}^{\infty}(-1)^k\left(\frac{z}{3}\right)^k=\sum_{k=0}^{\infty}\frac{(-1)^k}{3^{k+1}}z^k$$

收斂條件：$\left|\dfrac{z}{3}\right|<1$，即 $|z|<3$。

奇異點 $z=-3$，距原點 $=3$。
\end{solbox}

\begin{ansbox}
$R=\boxed{3}$
\end{ansbox}

%==================================================================
% Q6
%==================================================================
\begin{qbox}[第 6 題\,---\,圍道積分與奇異點]
$C:|z-\ii|=1$，計算 $\displaystyle\oint_C\frac{\ee^z}{z^2+1}\dd{z}=k\pi$，求 $k$。
\end{qbox}

\begin{solbox}[解]
\textbf{Step 1：找奇異點}

$\dfrac{\ee^z}{z^2+1}=\dfrac{\ee^z}{(z-\ii)(z+\ii)}$，奇異點為 $z=\ii$ 和 $z=-\ii$（皆為一階極點）。

\textbf{Step 2：判斷哪些在 $C:|z-\ii|=1$ 內}

\begin{itemize}
\item $z=\ii$：$|\ii-\ii|=0<1$ $\Rightarrow$ \textbf{在圓內}
\item $z=-\ii$：$|-\ii-\ii|=|-2\ii|=2>1$ $\Rightarrow$ 在圓外
\end{itemize}

\textbf{Step 3：利用 Cauchy 積分公式（Ch 5.5，Theorem 5.5.1）}

將被積函數寫成 $\dfrac{g(z)}{z-\ii}$ 的形式，其中 $g(z)=\dfrac{\ee^z}{z+\ii}$ 在 $C$ 內解析：

$$\oint_C\frac{\ee^z}{(z+\ii)(z-\ii)}\dd{z}=\oint_C\frac{g(z)}{z-\ii}\dd{z}=2\pi\ii\cdot g(\ii)$$

\textbf{Step 4：計算 $g(\ii)$}

$$g(\ii)=\frac{\ee^\ii}{\ii+\ii}=\frac{\ee^\ii}{2\ii}$$

\textbf{Step 5：代入}

$$\oint_C\frac{\ee^z}{z^2+1}\dd{z}=2\pi\ii\cdot\frac{\ee^\ii}{2\ii}=\pi\ee^\ii=\pi(\cos 1+\ii\sin 1)$$

展開為實部和虛部：$\pi\cos 1+\pi\ii\sin 1\approx 1.696+2.643\ii$。
\end{solbox}

\begin{ansbox}
$\displaystyle\oint_C\frac{\ee^z}{z^2+1}\dd{z}=\boxed{\pi\ee^\ii=\pi\cos 1+\pi\ii\sin 1}$
\end{ansbox}

%==================================================================
% Q7
%==================================================================
\begin{qbox}[第 7 題\,---\,殘值]
求 $\ee^{2/z}$ 在 $z=0$ 的殘值。
\end{qbox}

\begin{solbox}[解]
\textbf{Step 1：判斷奇異點類型}

$\ee^{2/z}$ 在 $z=0$ 非解析。利用 $\ee^w$ 的 Maclaurin 展開（Ch 6.2），令 $w=2/z$ 展開成 Laurent 級數。

\textbf{Step 2：Laurent 級數展開}

$\ee^w=\displaystyle\sum_{k=0}^{\infty}\frac{w^k}{k!}$（對所有 $w$ 收斂），代入 $w=\dfrac{2}{z}$：

$$\ee^{2/z}=\sum_{k=0}^{\infty}\frac{1}{k!}\left(\frac{2}{z}\right)^k=\underbrace{1}_{k=0}+\underbrace{\frac{2}{z}}_{k=1}+\underbrace{\frac{4}{2!\,z^2}}_{k=2}+\underbrace{\frac{8}{3!\,z^3}}_{k=3}+\cdots$$

$$=1+\frac{2}{z}+\frac{2}{z^2}+\frac{4}{3z^3}+\cdots$$

負冪次項有無窮多個，故 $z=0$ 為\textbf{本質奇異點}（Ch 6.4）。

\textbf{Step 3：讀取殘值}

殘值 $=a_{-1}=$ Laurent 級數中 $\dfrac{1}{z}$ 的係數 $=\dfrac{2^1}{1!}=2$。
\end{solbox}

\begin{ansbox}
$\Res(\ee^{2/z},0)=\boxed{2}$
\end{ansbox}

%==================================================================
% Q8
%==================================================================
\begin{qbox}[第 8 題\,---\,Laurent 展開]
將 $f(z)=\dfrac{1}{z(z-2)}$ 以 $|z|>2$ 展開成 Laurent series，求第四項。
\end{qbox}

\begin{solbox}[解]
$$\frac{1}{z(z-2)}=\frac{1}{z}\cdot\frac{1}{z-2}=\frac{1}{z^2}\cdot\frac{1}{1-2/z}=\frac{1}{z^2}\sum_{k=0}^{\infty}\left(\frac{2}{z}\right)^k=\sum_{k=0}^{\infty}\frac{2^k}{z^{k+2}}$$

$$=\frac{1}{z^2}+\frac{2}{z^3}+\frac{4}{z^4}+\frac{8}{z^5}+\cdots$$

第四項（$k=3$）：$\dfrac{8}{z^5}$。
\end{solbox}

\begin{ansbox}
第四項 $=\boxed{\dfrac{8}{z^5}}$
\end{ansbox}

%==================================================================
% Q9
%==================================================================
\begin{qbox}[第 9 題\,---\,實數積分]
$\displaystyle\int_0^{\infty}\frac{3\sin x}{x}\dd{x}=?$
\end{qbox}

\begin{solbox}[解]
\textbf{Step 1：利用 Dirichlet 積分（可由 Ch 6.6 的 indented contour 方法推導）}

考慮 $f(z)=\dfrac{\ee^{\ii z}}{z}$，利用 indented contour（實軸 $[-R,-r]\cup[r,R]$、大上半圓 $C_R$、小上半圓 $C_r$），並結合：
\begin{itemize}
\item Cauchy-Goursat 定理（$z=0$ 排除在圍道外，圍道內無奇異點）
\item Theorem 6.6.2：$R\to\infty$ 時大半圓積分 $\to 0$
\item Theorem 6.6.3：$r\to 0$ 時小半圓積分 $\to -\pi\ii\cdot\Res(f,0)=-\pi\ii$
\item 實軸積分合併：$\displaystyle\int_r^R\frac{\ee^{\ii x}-\ee^{-\ii x}}{x}\dd{x}=\int_r^R\frac{2\ii\sin x}{x}\dd{x}$
\end{itemize}

取極限得：$2\ii\displaystyle\int_0^{\infty}\frac{\sin x}{x}\dd{x}-\pi\ii=0$，即

$$\int_0^{\infty}\frac{\sin x}{x}\dd{x}=\frac{\pi}{2}$$

（完整推導見考古題 Q9 詳解。）

\textbf{Step 2：代入原題}

$$\int_0^{\infty}\frac{3\sin x}{x}\dd{x}=3\cdot\frac{\pi}{2}=\frac{3\pi}{2}$$
\end{solbox}

\begin{ansbox}
$\boxed{\dfrac{3\pi}{2}\approx 4.712}$
\end{ansbox}

%==================================================================
% Q10
%==================================================================
\begin{qbox}[第 10 題\,---\,三角積分]
$\displaystyle\int_0^{2\pi}\frac{1}{5+4\cos\theta}\dd{\theta}=?$
\end{qbox}

\begin{solbox}[解]
\textbf{Step 1：$z=\ee^{\ii\theta}$ 代換（Ch 6.6.1 方法）}

令 $z=\ee^{\ii\theta}$，則 $\cos\theta=\dfrac{z+z^{-1}}{2}$，$\dd{\theta}=\dfrac{\dd{z}}{\ii z}$，$C:|z|=1$。

\textbf{Step 2：轉換為圍道積分}

$$\int_0^{2\pi}\frac{1}{5+4\cos\theta}\dd{\theta}=\oint_{|z|=1}\frac{1}{5+4\cdot\frac{z+z^{-1}}{2}}\cdot\frac{\dd{z}}{\ii z}=\oint_{|z|=1}\frac{1}{5+2z+2z^{-1}}\cdot\frac{\dd{z}}{\ii z}$$

\textbf{Step 3：化簡分母（消去負冪次）}

$\ii z(5+2z+2z^{-1})=\ii(5z+2z^2+2)=2\ii(z^2+\tfrac{5}{2}z+1)$

故原式 $=\displaystyle\oint_{|z|=1}\frac{\dd{z}}{2\ii(z^2+\frac{5}{2}z+1)}=\frac{1}{2\ii}\oint_{|z|=1}\frac{\dd{z}}{z^2+\frac{5}{2}z+1}$

\textbf{Step 4：因式分解找極點}

$z^2+\dfrac{5}{2}z+1=0 \Rightarrow z=\dfrac{-5/2\pm\sqrt{25/4-4}}{2}=\dfrac{-5/2\pm 3/2}{2}$

\begin{itemize}
\item $z=\dfrac{-5/2+3/2}{2}=\dfrac{-1}{2}$：$\left|-\dfrac{1}{2}\right|=\dfrac{1}{2}<1$ $\Rightarrow$ \textbf{在單位圓內}
\item $z=\dfrac{-5/2-3/2}{2}=-2$：$|-2|=2>1$ $\Rightarrow$ 在單位圓外
\end{itemize}

故 $z^2+\dfrac{5}{2}z+1=(z+\tfrac{1}{2})(z+2)$。

\textbf{Step 5：計算殘值}

$z=-\dfrac{1}{2}$ 為一階極點，由 Theorem 6.5.1：
$$\Res\!\left(\frac{1}{(z+\frac{1}{2})(z+2)},\,-\frac{1}{2}\right)=\lim_{z\to -1/2}\frac{(z+\frac{1}{2})}{(z+\frac{1}{2})(z+2)}=\frac{1}{-\frac{1}{2}+2}=\frac{1}{\frac{3}{2}}=\frac{2}{3}$$

\textbf{Step 6：殘值定理求積分}

$$\frac{1}{2\ii}\oint_{|z|=1}\frac{\dd{z}}{(z+\frac{1}{2})(z+2)}=\frac{1}{2\ii}\cdot 2\pi\ii\cdot\frac{2}{3}=\frac{2\pi}{3}$$
\end{solbox}

\begin{ansbox}
$\boxed{\dfrac{2\pi}{3}}$
\end{ansbox}

%==================================================================
% Q11
%==================================================================
\begin{qbox}[第 11 題\,---\,線性分式變換]
$|z|=1$ 在 $T(z)=\dfrac{z+1}{z-1}$ 映射後為？
\end{qbox}

\begin{solbox}[解]
令 $z=\ee^{\ii\theta}$，$|z|=1$：

$$w=\frac{\ee^{\ii\theta}+1}{\ee^{\ii\theta}-1}=\frac{(\cos\theta+1)+\ii\sin\theta}{(\cos\theta-1)+\ii\sin\theta}$$

分子分母乘以共軛 $(\cos\theta-1)-\ii\sin\theta$：

分子 $=(\cos\theta+1)(\cos\theta-1)-\ii(\cos\theta+1)\sin\theta+\ii\sin\theta(\cos\theta-1)+\sin^2\theta$
$=\cos^2\theta-1+\sin^2\theta+\ii\sin\theta[(\cos\theta-1)-(\cos\theta+1)]$
$=0+\ii(-2\sin\theta)=-2\ii\sin\theta$

分母 $=(\cos\theta-1)^2+\sin^2\theta=2-2\cos\theta$

$$w=\frac{-2\ii\sin\theta}{2(1-\cos\theta)}=\frac{-\ii\sin\theta}{1-\cos\theta}$$

$w$ 為純虛數，故映射為\textbf{虛軸}。
\end{solbox}

\begin{ansbox}
映射後為\boxed{\text{虛軸（}\text{Re}(w)=0\text{）}}
\end{ansbox}

%==================================================================
% Q12
%==================================================================
\begin{qbox}[第 12 題\,---\,8 字形圍道]
$C$ 為 8 字形，$C_1$（逆時針包含 $z=0$）和 $C_2$（順時針包含 $z=2$）。

計算 $\displaystyle\oint_{C_1}\frac{3z-4}{z^2-2z}\dd{z}=k_1\cdot 2\pi\ii$ 和 $\displaystyle\oint_{C_2}\frac{3z-4}{z^2-2z}\dd{z}=k_2\cdot 2\pi\ii$。
\end{qbox}

\begin{solbox}[解]
$\dfrac{3z-4}{z^2-2z}=\dfrac{3z-4}{z(z-2)}=\dfrac{A}{z}+\dfrac{B}{z-2}$

$A=\dfrac{3(0)-4}{0-2}=2$，$B=\dfrac{3(2)-4}{2}=1$。

\textbf{$C_1$}（逆時針，含 $z=0$）：
$$\oint_{C_1}\frac{3z-4}{z(z-2)}\dd{z}=2\cdot 2\pi\ii+0=4\pi\ii$$
$k_1=2$。

\textbf{$C_2$}（順時針，含 $z=2$）：
$$\oint_{C_2}\frac{3z-4}{z(z-2)}\dd{z}=0+1\cdot(-2\pi\ii)=-2\pi\ii$$
$k_2=-1$。
\end{solbox}

\begin{ansbox}
$k_1=\boxed{2}$，$k_2=\boxed{-1}$
\end{ansbox}

%==================================================================
% Q13
%==================================================================
\begin{qbox}[第 13 題\,---\,混合積分]
$C:|z|=1$，計算 $\displaystyle\oint_C\left(\frac{3}{z}+\cos z+z^2+1\right)\dd{z}=k\cdot 2\pi\ii$，求 $k$。
\end{qbox}

\begin{solbox}[解]
由積分的線性性質，逐項計算：

\textbf{(i)} $\displaystyle\oint_C\frac{3}{z}\dd{z}$：$z=0$ 在 $|z|=1$ 內部，為 $\dfrac{3}{z}$ 的一階極點。

由 $\oint_C\dfrac{1}{z}\dd{z}=2\pi\ii$（Ch 5.3），得 $\oint_C\dfrac{3}{z}\dd{z}=3\cdot 2\pi\ii=6\pi\ii$。

\textbf{(ii)} $\displaystyle\oint_C\cos z\dd{z}=0$：$\cos z$ 為 entire function（在整個複數平面解析），由 Cauchy-Goursat 定理（Ch 5.3），解析函數沿封閉曲線的積分為零。

\textbf{(iii)} $\displaystyle\oint_C z^2\dd{z}=0$：$z^2$ 為 entire function，同理。

\textbf{(iv)} $\displaystyle\oint_C 1\dd{z}=0$：常數函數亦為 entire function，同理。

合併：
$$\oint_C\left(\frac{3}{z}+\cos z+z^2+1\right)\dd{z}=6\pi\ii+0+0+0=6\pi\ii=3\cdot 2\pi\ii$$
\end{solbox}

\begin{ansbox}
$k=\boxed{3}$
\end{ansbox}

%==================================================================
% Q14
%==================================================================
\begin{qbox}[第 14 題\,---\,殘值計算]
(a) $f(z)=\dfrac{3z+5}{z-2}$，求 $\Res(f(z),2)$。

(b) $f(z)=(z-1)^2\cos\!\left(\dfrac{1}{z-1}\right)$，求 $\Res(f(z),1)$。
\end{qbox}

\begin{solbox}[解]
\textbf{(a)} $f(z)=\dfrac{3z+5}{z-2}$，$z=2$ 為一階極點。

由 Theorem 6.5.1（一階極點殘值公式）：
$$\Res(f(z),2)=\lim_{z\to 2}(z-2)\cdot\frac{3z+5}{z-2}=\lim_{z\to 2}(3z+5)=3(2)+5=11$$

\textbf{(b)} $f(z)=(z-1)^2\cos\!\left(\dfrac{1}{z-1}\right)$，$z=1$ 為奇異點。

\textbf{Step 1}：令 $w=z-1$（平移至 $w=0$），則 $f=w^2\cos\!\left(\dfrac{1}{w}\right)$。

\textbf{Step 2}：利用 $\cos u$ 的 Maclaurin 展開（Ch 6.2），令 $u=\dfrac{1}{w}$：
$$\cos\!\left(\frac{1}{w}\right)=1-\frac{1}{2!\,w^2}+\frac{1}{4!\,w^4}-\frac{1}{6!\,w^6}+\cdots$$

\textbf{Step 3}：乘以 $w^2$ 得 Laurent 級數：
$$w^2\cos\!\left(\frac{1}{w}\right)=w^2\cdot 1-\frac{w^2}{2!\,w^2}+\frac{w^2}{4!\,w^4}-\frac{w^2}{6!\,w^6}+\cdots$$
$$=w^2-\frac{1}{2}+\frac{1}{24w^2}-\frac{1}{720w^4}+\cdots$$

\textbf{Step 4}：觀察 Laurent 級數，\textbf{沒有 $\dfrac{1}{w}$ 項}（$w^{-1}$ 的係數為 $0$）。

這是因為 $\cos$ 的 Maclaurin 展開只含\textbf{偶數}冪次項，乘以 $w^2$ 後只產生偶數冪次，不會出現奇數負冪次如 $w^{-1}$。
\end{solbox}

\begin{ansbox}
(a) $\Res(f(z),2)=\boxed{11}$

(b) $\Res(f(z),1)=\boxed{0}$
\end{ansbox}

\end{document}
